% !PRINTED 344
\weberpage{344}

Finally we form one more triple containing $(0,\,2)$ and one
containing $(2,\,0)$,
\[
(0,\,2)\ (1,\,2)\ (2,\,2), \quad (2,\,0)\ (2,\,1)\ (2,\,2),
\]
with which all twelve triples, and with them the arrangement \S.\,91, (20),
are determined.

\S.\,93.

\gesperrt{The group of the triple equations.}

We shall now determine the \gesperrt{Galois} group $P$ of the triple
equations. First of all it is clear that in this group only
such permutations occur under which every triple passes again
into a triple. For if we assume that the roots of a
triple $x_1, x_2, x_3$ pass by a permutation into $y_1, y_2, y_3$,
then, since this permutation is applicable to the equation (1) (\S.\,92),
it follows that
\[
y_1 = \Theta\,(y_2, y_3),
\]
and consequently $y_1$ is the third root of the triple determined
by $y_2, y_3$.

If we now pass over to the notation for the roots:
\weq{1}{x = (\xi, \eta)}
then it results from the compilation of the triples in
§§.\,91 and 92 that three roots
\weq{2}{x_1 = (\xi_1, \eta_1), \quad x_2 = (\xi_2, \eta_2), \quad x_3 = (\xi_3, \eta_3)}
always form a triple then, and only then, when
\weq{3}{\begin{array}{l}
\xi_1 + \xi_2 + \xi_3 \equiv 0\\
\eta_1 + \eta_2 + \eta_3 \equiv 0
\end{array} \quad (\mathrm{mod}\ 3).}
If we call a permutation group of the $(\xi, \eta)$ whose substitutions
all belong to the complete linear congruence group for
the modulus 3 (\S.\,78) a \gesperrt{linear group}, then the following theorem
holds:

\gesperrt{1. The group $P$ of a triple equation is always
linear.}

According to \S.\,77, every permutation of the quantities $(\xi, \eta)$
whatever can be represented in the form:
\weq{4}{\xi' \equiv \varphi\,(\xi, \eta), \quad \eta' \equiv \psi\,(\xi, \eta) \quad (\mathrm{mod}\ 3),}
%CONT>>
